
% Exact finite vacuum support and adelic cone factorisation.
% Scope: closes the analytic/arithmetic content of ConeFacA and ConeFacB,
% given the explicitly constructed evaluation and oscillator models already
% established in the preceding sections. It does not by itself identify
% global theta constant terms or prove OneConst outside the compact base.

\section{The coefficient ideal and the adelic cone integral}
\label{sec:sr-cone-euler}

We use the oscillator, evaluation, and compact stabilizer constructed above.
Let \(E/F\) be the CM extension of the fixed base. In particular it is
unramified at every finite place. Write
\(\mathfrak c=\mathfrak d_{E/\mathbb Q}^{-1}\).
An adapted hyperbolic plane in a Hermitian space \(X\) has basis \(x,y\),
with \(\langle x,y\rangle=1\), and coefficient ideal \(\mathfrak b\)
on \(x\). An adapted plane in the skew Hermitian space \(W\) has basis
\(x',y'\), with \(\langle x',y'\rangle=1\), and coefficient ideal
\(\mathfrak b'\) on \(x'\). Let \(X_1,W_1\) be the specified orthogonal
complements. All the lattice decompositions here are the adapted
decompositions, with their specified embeddings into the upper spaces.

The two shifts, in the same tensor symplectic space and the same genuine
oscillator representation as the evaluation operator, are
\[
 u_A(l)=l\,x\otimes y',\qquad
 u_B(l)=\bar l\,y\otimes x'.
\]
The respective evaluation isotropic spaces are
\[
 I_A=X\otimes Ex'+Ex\otimes W_1,\qquad
 I_B=Ex\otimes W+X_1\otimes Ex'.
\]
Their symplectic reductions are the specified \(X_1\otimes W_1\).
The local evaluation maps are normalized to send the upper lattice vacuum
to the lower lattice vacuum. This normalization is used at every finite
place, including those above \(2\).

\begin{srlemma}[The two coefficient ideals]
\label{lem:sr-cone-coefficient-ideals}
Put
\[
 \mathfrak a_A=\mathfrak b\,\mathfrak c\,
                     \bar{\mathfrak b'}^{-1},
 \qquad
 \mathfrak a_B=
       \overline{\bar{\mathfrak b}^{-1}\mathfrak b'}
       =\mathfrak b^{-1}\bar{\mathfrak b'}.
\]
These are nonzero fractional ideals. At every finite place \(v\) of \(F\),
\[
 u_\epsilon(l)_v\in(\Lambda_X\otimes_{\mathcal O_E}\Lambda_W)_v
 \quad\Longleftrightarrow\quad
 l_w\in(\mathfrak a_\epsilon)_w
       \quad\hbox{for every \(w\mid v\)},\qquad \epsilon=A,B.
 \tag{C1}\label{eq:sr-cone-coefficient-support}
\]
At the distinguished split places of the fixed base the valuations of
both ideals are zero.

Let \(\nu\) be a unitary Hecke character, trivial on
\(\mathbb A_F^\times\), unramified outside the distinguished places.
For an ideal with zero valuations at those places, define
\[
 \nu(\mathfrak a)=
 \prod_w\nu_w(\varpi_w)^{\operatorname{ord}_w\mathfrak a};
\]
this finite product is independent of the uniformizers. Then
\[
 \nu(\mathfrak a_A)=\nu(\mathfrak b)\nu(\mathfrak b'),
 \qquad
 \nu(\mathfrak a_B)=
             \nu(\mathfrak b)^{-1}\nu(\mathfrak b')^{-1}.
 \tag{C2}\label{eq:sr-cone-ideal-characters}
\]
\end{srlemma}

\begin{proof}
The duality conditions on the actual lattices give
\[
 \Lambda_X=
   \mathfrak b x\ \oplus\ \Lambda_{X_1}\
          \oplus\bar{\mathfrak b}^{-1}y,\qquad
 \Lambda_W=
   \mathfrak b'x'\ \oplus\ \Lambda_{W_1}\
          \oplus\mathfrak c\,\bar{\mathfrak b'}^{-1}y'.
 \tag{C3}\label{eq:sr-cone-lattice-splitting}
\]
For clarity, the last coefficient ideal in either formula is forced,
rather than chosen independently: pairing \(t y\) with all \(s x\),
\(s\in\mathfrak b\), requires
\(t\bar{\mathfrak b}\subset\mathcal O_E\), and the corresponding
trace duality for \(W\) requires
\(t\bar{\mathfrak b'}\subset\mathfrak c\).
Orthogonality and the adapted decomposition give the converse inclusions.

Tensoring these direct decompositions identifies the coefficient ideal
of \(x\otimes y'\) as
\(\mathfrak b\mathfrak c\bar{\mathfrak b'}^{-1}\), and that of
\(y\otimes x'\) as \(\bar{\mathfrak b}^{-1}\mathfrak b'\).
Conjugating the coefficient in the second shift gives (C1).
This argument can also be checked without a tensor decomposition:
choose integral coefficient projections onto each line. Applying the
tensor product of these projections gives necessity; products of elements
of the two coefficient ideals generate their ideal product and give
sufficiency. It works for fractional, nonprincipal ideals.

At a distinguished split place the adapted line vectors are columns of
the specified integral basis. Their coefficient projections are integral
and take value \(1\) on these columns, so the relevant line ideals have
valuation zero. The standard additive character there has integral
conductor, so the codifferent has valuation zero there as well.

The restriction of \(\nu\) gives
\(\nu(\bar{\mathfrak j})=\nu(\mathfrak j)^{-1}\) for every ideal under
consideration: a local generator times its conjugate comes from
\(F_v^\times\). Finally finite unramifiedness of \(E/F\) gives
\[
 \mathfrak d_{E/\mathbb Q}^{-1}
   =\mathfrak d_{F/\mathbb Q}^{-1}\mathcal O_E.
\]
Generators of the latter can be chosen in \(F_v^\times\), and therefore
\(\nu(\mathfrak c)=1\). Multiplicativity now gives (C2).
\end{proof}

\begin{srlemma}[Exact translated vacuum and finite cone operator]
\label{lem:sr-cone-finite-vacuum}
Let \(e_v\) be either of the normalized finite evaluation maps above and
let \(\phi_v^0,\phi_{1,v}^0\) be the actual upper and lower vacua. Then
\[
 e_v\rho_v(u_\epsilon(l))\phi_v^0
       =\mathbf1_{(\mathfrak a_\epsilon)_v}(l)\,\phi_{1,v}^0.
 \tag{C4}\label{eq:sr-cone-exact-vacuum}
\]
Here the indicator imposes its condition at all \(w\mid v\), and the
Heisenberg operator and the shift are in the same upper model.

Let \(K_{X,v}\) be the specified lattice stabilizer. Give it probability
Haar measure, and give each \(E_w^\times\) multiplicative Haar measure
with \(\operatorname{vol}\mathcal O_w^\times=1\).
For every integer \(z\ge2\), the formula
\[
 \Lambda_{\epsilon,v}\phi=
  \int_{K_{X,v}}\int_{(E\otimes_F F_v)^\times}
       e_v\rho_v(u_\epsilon(l))\omega_v(k,1)\phi\,
                  \nu_v(l)|l|_{E_v}^{z}\,d^\times l\,dk
 \tag{C5}\label{eq:sr-cone-finite-operator}
\]
defines a linear map of Schwartz--Bruhat spaces.
Whenever the two splitting characters and \(\nu\) are unramified at \(v\),
\[
 \Lambda_{\epsilon,v}\phi_v^0
     =u_{\epsilon,v}\phi_{1,v}^0,\quad
 u_{\epsilon,v}=
 \prod_{w\mid v}
 \frac{\bigl(\nu_w(\varpi_w)Nw^{-z}\bigr)^{
                       \operatorname{ord}_w\mathfrak a_\epsilon}}
      {1-\nu_w(\varpi_w)Nw^{-z}}
 \ne0.
 \tag{C6}\label{eq:sr-cone-finite-euler}
\]
\end{srlemma}

\begin{proof}
Work first at one \(w\). Localizing (C3) gives free orthogonal
decompositions, since \(\mathcal O_w\) is a discrete valuation ring.
In the partial Fourier coordinates defining \(e_v\), its eliminated
coordinates and their duals are complementary isotropic coordinates.
The corner \(x\otimes y'\), respectively \(y\otimes x'\), is one
of the dual coordinates; its dual coefficient lattice is exactly the
corresponding ideal in (C1). The upper vacuum is the characteristic
function of the upper polarized lattice, up to its fixed model scalar.
Integration in an eliminated coordinate has Fourier transform equal to
its volume times the characteristic function of the dual coefficient
lattice. Evaluation in the other eliminated coordinates gives their
values at zero. Consequently the translated evaluation is zero outside
(C1), and is a scalar multiple of the lower vacuum inside (C1).
Its scalar inside is the same as at \(l=0\): translation in a lattice
coordinate preserves its characteristic function and modulation on the
dual lattice is \(1\). The normalization
\(e_v\phi_v^0=\phi_{1,v}^0\) makes that scalar \(1\).
The corner has zero symplectic self-pairing, so no quadratic Heisenberg
phase occurs. This proves (C4), also in residue characteristic \(2\);
the dual lattice uses the actual trace character and codifferent and
does not require division of an integral coordinate by \(2\).

Here is also a direct finiteness argument for (C5) on arbitrary
Schwartz--Bruhat functions. A given function is fixed by Heisenberg
translation and modulation from a sufficiently small lattice.
Commuting such a modulation through the corner translation shows that
a nonzero evaluation requires \(l\) to lie in a fixed fractional ideal:
otherwise one of the invariant modulation characters is nontrivial.
On this bounded set, the vector valued map
\[
 l\longmapsto e_v\rho_v(u_\epsilon(l))\phi
\]
is locally constant, extends to \(l=0\), and has finite image in the
lower Schwartz--Bruhat space. For \(|l|\) sufficiently small it is
constant. Thus its integral is a finite linear combination of lower
Schwartz--Bruhat functions. The tail toward zero is absolutely
integrable, since on a shell of valuation \(j\) its absolute weight is
\(Nw^{-jz}\), and \(z>0\). The compact group \(K_{X,v}\) has finite orbit
on a fixed smooth vector, so its average preserves this conclusion.
These arguments also prove absolute integrability at each lower
coordinate, and justify the linearity of the integral.

At an unramified place the specified oscillator fixes \(\phi_v^0\)
under \(K_{X,v}\); this includes the prescribed unramified splitting
scalar, so it is equality of vectors. Insert (C4) into (C5). At \(w\)
the resulting integral is
\[
 \sum_{j\ge m_w}\nu_w(\varpi_w)^jNw^{-jz}
    =\frac{(\nu_w(\varpi_w)Nw^{-z})^{m_w}}
                  {1-\nu_w(\varpi_w)Nw^{-z}},\qquad
       m_w=\operatorname{ord}_w\mathfrak a_\epsilon.
\]
The absolute value of the denominator parameter is \(Nw^{-z}<1\).
Multiplying over \(w\mid v\) proves (C6).
\end{proof}

We now use the archimedean cone map already constructed by partial
Fourier transformation, Gaussian moments, and compact averaging:
\[
 L_{\epsilon,\infty}f=
  \int_{K_{X,\infty}}\int_{E_\infty^\times}
 e_\infty\rho_\infty(u_\epsilon(l))\omega_\infty(k,1)f\,
             \nu_\infty(l)|l|_{E_\infty}^{z}\,d^\times l\,dk.
 \tag{C7}\label{eq:sr-cone-arch-operator-again}
\]
It is a continuous map of Schwartz spaces and preserves the specified
polynomial Gaussian subspaces. Its defining integrals converge in the
Schwartz seminorms. The strict positivity of the exponent at zero and
Schwartz decay along the nonzero corner give these facts, as proved
above; no assertion about a global theta integral is used here.

\begin{srproposition}[Adelic factorization with constants before packets]
\label{prop:sr-cone-euler-factorization}
Fix the base, the two adapted planes, the actual upper and lower models,
their evaluation map \(e\), and \(z\ge2\). Fix Haar measures \(\mu_K\)
on the specified compact stabilizer \(K_X\) and \(\mu_\times\) on
\(\mathbb A_E^\times\), and fix the local measures used in (C5) and
(C7). There is a nonzero scalar \(\kappa_\epsilon\), depending on these
fixed data and \(z\), before choosing the finite character packet, with
the following property.

Let \(\nu\) be as in the coefficient ideal lemma. Let \(S_0\subset S\)
be finite sets of finite places of \(F\) such that \(\nu\) and the
splitting characters are unramified off \(S\), all the local factors
of \(\nu\) above \(S\) are ramified, and
\(\operatorname{ord}_w\mathfrak a_\epsilon=0\) for \(w\) above \(S\).
For any finite \(U\supset S\) and pure Schwartz vector
\(\Phi=f\otimes\bigotimes_v\phi_v\), with
\(\phi_v=\phi_v^0\) off \(U\), put
\[
 \begin{split}
 C_\epsilon\Phi={}&
 \int_{K_X}\int_{\mathbb A_E^\times}
 e\rho(u_\epsilon(l))\omega(k,1)\Phi\,
                  \nu(l)|l|_E^z\,d\mu_\times(l)\,d\mu_K(k),\\
 \Psi_U={}& L_{\epsilon,\infty}f\ \otimes
        \bigotimes_{v\in U}\Lambda_{\epsilon,v}\phi_v
              \ \otimes\bigotimes_{v\notin U}\phi_{1,v}^0.
 \end{split}
\]
The integrals defining \(C_\epsilon\Phi\) converge, define a Schwartz
vector, and satisfy
\[
 \left(\prod_{v\in U\setminus S}u_{\epsilon,v}\right)
                 C_\epsilon\Phi
   =\kappa_\epsilon\,\nu(\mathfrak a_\epsilon)\,
                      L^{S_0}(\nu,z)\,\Psi_U.
 \tag{C8}\label{eq:sr-cone-global-factorization}
\]
They extend linearly to the adelic Schwartz space. In particular their
coefficient ideal factors are exactly
\[
 \nu(\mathfrak b)\nu(\mathfrak b')\quad\hbox{for \(A\)},\qquad
 [\nu(\mathfrak b)\nu(\mathfrak b')]^{-1}
                  \quad\hbox{for \(B\)}.
 \tag{C9}\label{eq:sr-cone-fixed-ideal-exponents}
\]
The maps \(L_{\epsilon,\infty}\), \(\Lambda_{\epsilon,v}\), and the
numbers \(u_{\epsilon,v}\) are the actual integral operators and actual
vacuum factors of (C5)--(C7), rather than separately chosen matching
maps.
\end{srproposition}

\begin{proof}
The compact stabilizer condition is a condition at each completion:
the fixed positive Hermitian metric at infinity and the specified
lattice at every finite place. It therefore identifies \(K_X\)
with the product of these compact local stabilizers. Its probability
Haar measure is their product; comparison with \(\mu_K\) contributes a
positive scalar depending only on the fixed Haar measures.
Likewise the product of the chosen infinite multiplicative measure and
the finite measures of unit volume defines Haar measure on the
restricted product \(\mathbb A_E^\times\). Comparison with
\(\mu_\times\) contributes a second positive scalar, also independent
of every character and test vector. Denote their product by \(c_H>0\).

Oscillator, Heisenberg, and evaluation actions factor in these product
coordinates: their operators are the genuine ones already normalized
in a common tensor symplectic model. Thus, off \(U\), their integrand is
exactly the lower vacuum times the indicator in (C4). The finite
operator lemma treats the places in \(U\), and (C7) treats infinity.

There is an absolute majorant for the remaining restricted product
integral. If \(m_w=\operatorname{ord}_w\mathfrak a_\epsilon\), its
integral at \(w\) is
\[
  \sum_{j\ge m_w}Nw^{-jz}
          =Nw^{-m_wz}(1-Nw^{-z})^{-1}.
\]
Only finitely many \(m_w\) are nonzero. The product of the other factors
converges for \(z\ge2\): the prime ideal sum
\(\sum_wNw^{-z}\) is finite, as follows by bounding the logarithm of
the absolutely convergent Dedekind zeta Euler product there.
This gives an integrable majorant on the restricted product, multiplied
by the finite collection of integrable factors already considered.
Consequently exhaustion by finite sets of places, Fubini, and dominated
convergence all apply. They apply also to each lower Schwartz seminorm:
outside \(U\) the lower factor is a fixed vacuum; at the finitely many
remaining places we have the finite dimensional image in the proof of
(C5), and at infinity (C7) converges in every Schwartz seminorm.
Thus the result is a Schwartz vector.

The resulting equality, before multiplying any finite vacuum factors,
is
\[
 C_\epsilon\Phi
      =c_H\left(\prod_{v\notin U}u_{\epsilon,v}\right)\Psi_U.
 \tag{C10}\label{eq:sr-cone-before-euler}
\]
The product in (C10) includes the complete geometric series integrals;
it is not an assertion that an unramified integral has compact support.
The absolutely convergent Euler calculation (C6) gives
\[
 \prod_{v\notin S}u_{\epsilon,v}
   =\left(\prod_wNw^{-zm_w}\right)
        \left(\prod_w\nu_w(\varpi_w)^{m_w}\right)
                         L^S(\nu,z)
   =N(\mathfrak a_\epsilon)^{-z}
             \nu(\mathfrak a_\epsilon)L^S(\nu,z).
 \tag{C11}\label{eq:sr-cone-complete-euler}
\]
The zero valuation condition over \(S\) permits the products in the
middle expression to run over all \(w\). The ramified local Euler
factors over \(S\setminus S_0\) equal \(1\), so
\(L^S(\nu,z)=L^{S_0}(\nu,z)\).
Multiply (C10) by the finite product in (C8) and use (C11). This proves
(C8) with
\[
 \kappa_\epsilon=c_H N(\mathfrak a_\epsilon)^{-z}>0.
\]
The ideal, its norm, the exponent \(z\), and the Haar comparisons were
fixed before \(\nu\); none of their factors depends on the finite
character packet. Formula (C9) follows from (C2).

Finally each finite adelic Schwartz--Bruhat function is a finite sum of
characteristic functions of compact open cosets. Each such coset
splits into finitely many boxes at a finite set of places and the
standard lattice outside that set. Hence every adelic Schwartz vector
is a finite sum of the pure vectors just considered. The absolute
convergence already proved permits addition and scalar multiplication
under the integral and proves the stated extension.
\end{proof}

\paragraph{Exact scope for the rank recursions.}
\label{rem:sr-cone-euler-scope}
The powers in (C9) are universal: \((1,1)\) for the \(A\) cone and
\((-1,-1)\) for the \(B\) cone. The sign in each ramified standard-vector
Gauss factor remains the \(+1\) already calculated with the actual
additive character. This section supplies the adelic cone factorization
and its packet quantifiers. To deduce a global rank recursion one must
still identify the appropriate global theta constant-term component
with this cone integral, and prove the descent through the specified
upper matching quotient. These are separate mathematical statements;
neither follows solely from an Euler product or a one-dimensional
space of local intertwining maps.

